\section{Experiments}
\label{sec:experiments}
We ask whether the placement of exact calls, rather than the reconstruction of skipped steps, explains what a cache can achieve at a fixed budget, and whether a schedule derived from the sampler's sensitivity profile beats the schedules that existing caches use.

\paragraph{Setup.} Every arm runs inside one implementation, the released CogVideoX pipeline of RA-CFGCache with its own reference, so that a comparison never mixes two code bases. Our schedule enters that implementation through a single hook that replaces the decision of whether to evaluate, and leaves the pipeline, the guidance, the decoder and the skip semantics untouched. We report LPIPS with a VGG backbone against a full run of the same pipeline, count a step as one transformer call, and use the first 32 prompts of VBench for CogVideoX-2B and the first 16 for the other two settings.

\begin{figure}[t]
\centering
\includegraphics[width=\linewidth]{figures/fig_sensitivity.pdf}
\caption{(a) The cost of skipping exactly one step, measured against a full run, for two prompts: the cost is concentrated in the head and the two profiles agree on its order. (b) The steps each rule evaluates at a comparable budget; the threshold rule keeps four steps in the head, while the measured profile asks for twelve.}
\label{fig:sensitivity}
\end{figure}

\paragraph{Main result.} Table~\ref{tab:main} reports three settings on all three metrics. On CogVideoX-2B at 17 frames, over one hundred prompts, our schedule reaches 0.145 LPIPS, 0.906 SSIM and 28.6 PSNR against 0.159, 0.895 and 27.7 for the strongest baseline, the risk-based cache of RA-CFGCache; the threshold-based cache falls to 0.279 and the Taylor expansion to 0.492. At 49 frames over 48 prompts it reaches 0.144 against 0.155, and on Wan2.1 at 81 frames 0.134 against 0.154. The ordering is the same on every metric and every setting, which is what a placement effect should look like: it does not trade one axis of quality for another. The improvement over the strongest baseline is between $1.09$ and $1.15$ times in LPIPS. Figure~\ref{fig:curve} sweeps the budget rather than comparing a single operating point: on both backbones the curve stays below every released baseline at the budget those baselines use.

\begin{table*}[t]
\centering\footnotesize\setlength{\tabcolsep}{5pt}
\begin{tabular}{llcccc}
\toprule
Setting & Method & Calls & LPIPS $\downarrow$ & SSIM $\uparrow$ & PSNR $\uparrow$ \\
\midrule
Cog 17f (100) & TaylorSeer & 19 & 0.492 & 0.526 & 13.9 \\
              & TeaCache & 17 & 0.279 & 0.752 & 21.1 \\
              & RA-CFGCache & 19 & 0.159 & 0.895 & 27.7 \\
              & \method{} & 20 & \textbf{0.145} & \textbf{0.906} & \textbf{28.6} \\
\midrule
Cog 49f (48)  & RA-CFGCache & 19 & 0.155 & 0.904 & 29.1 \\
              & \method{} & 20 & \textbf{0.144} & \textbf{0.914} & \textbf{29.9} \\
\midrule
Wan 81f (16)  & TaylorSeer & 19 & 0.379 & 0.620 & 17.6 \\
              & TeaCache & 19 & 0.199 & 0.827 & 24.0 \\
              & RA-CFGCache & 19 & 0.154 & 0.870 & 27.8 \\
              & \method{} & 21 & \textbf{0.134} & \textbf{0.888} & \textbf{28.9} \\
\bottomrule
\end{tabular}
\caption{All arms inside one pipeline and one reference, at a matched number of transformer calls; the number in parentheses is the prompt count. Our schedule is first on all three metrics in every setting. LPIPS uses a VGG backbone.}
\label{tab:main}
\end{table*}

\begin{figure}[t]
\centering
\includegraphics[width=\linewidth]{figures/fig_curve2.pdf}
\caption{LPIPS against the number of transformer calls at 16 prompts. The red curve varies only the length of our schedule; the three markers are the released baselines at their own default budgets, which their public entry points do not expose as a continuous knob. At seventeen calls our schedule is 20\% closer to the full run than the threshold-based cache at the same cost on CogVideoX-2B, and at nineteen calls it is below the risk-based cache on both backbones.}
\label{fig:curve}
\end{figure}

\paragraph{The tail must grow, not spread.} At a fixed budget the only freedom is the shape of the gaps after the head. We compare a uniform tail with the growing tail of Eq.~\ref{eq:sched}, both with the same head and the same number of calls, and the growing tail lowers LPIPS by 34\% relative to a uniform tail on the same twelve prompts. \paragraph{The interval sweep, and what it does not show.} A natural control is the family of fixed-interval schedules: the Taylor expansion baseline with its interval set to two, three and four, which costs about 26, 19 and 15 calls. On 32 prompts it reaches 0.402, 0.539 and 0.615, so our schedule is ahead at every interval. That comparison is not evidence for the shape of the tail, because the degradation of that baseline is dominated by its forecast rather than by its spacing: even at 26 calls, where its spacing is denser than ours, it stays at 0.40. The shape is isolated by the control above, which holds the head and the budget fixed and varies only the gaps.

The threshold rule fails the same way for a different reason. Its first exact calls are driven by the change of the input, which is largest at the start, but its budget tightens after four steps, and from there it spaces the rest by a threshold rather than by the profile: Figure~\ref{fig:sensitivity}b shows the four-step head against the twelve the profile asks for. Keeping more of the head exact is what the growing-gap rule buys.

\paragraph{The head length follows from the profile.} We measure the single-skip profile once per backbone and read the head length from it, without search: the first ten steps of CogVideoX-2B carry about three quarters of the total sensitivity, and the schedule therefore holds them exact. With a budget below the head length the rule keeps the earliest steps exact and accepts long gaps later, which is the regime where the gap between schedules is largest: at ten calls the schedule is 5.0 times closer to the full run than a uniform grid of the same size.

\paragraph{Scope.} The measurement and the schedule are verified on two backbones and two frame counts, with budgets between ten and twenty-one calls. A comparison at a much lower budget is not covered, and the profile is measured for the specific step count of each setting, so a different scheduler requires re-measuring it once. The absolute values of the baselines in Table~\ref{tab:main} come from our reproduction inside one harness and are not the numbers those papers report.
